% ============================================================
%  Chapter 7 — The Hidden Sensors
%  Inside a Smartphone: The Tiny World in Your Hand
% ============================================================

\chapter{The Hidden Sensors}

\chapteropening{%
  Your phone can feel you turning it sideways. It dims its screen
  in your pocket. It counts your steps without you pressing anything.
  Meet the invisible helpers doing all of this — the sensors.
}
\chapterbadge{Mission: Feel the World Through Your Phone!}

% --- Opening ---
\section*{Meet Sensa}

\character{Sensa}{I am Sensa — actually, I am a whole family of tiny sensors!
  Most people never notice us, but we are always watching and listening.
  We tell the phone which way is up, how fast it is moving,
  how bright the room is, and which direction is north.
  Without us, your phone would be blind to the physical world.}

\sectionrule

% --- Accelerometer ---
\section*{Sensor 1: The Accelerometer}

The \term{accelerometer} measures how fast the phone is speeding up or slowing down
— and in which direction.
Because gravity always pulls downward, it also helps the phone work out which way is up.

\begin{itemize}
  \item It detects when you rotate the phone from portrait to landscape
        and tells the screen to rotate too.
  \item It counts your steps in a fitness app by detecting each footfall.
  \item It detects a sudden drop and can take protective action.
\end{itemize}

\begin{experiment}
  Hold a phone and slowly rotate it from portrait (upright) to landscape (sideways).
  Watch the screen rotate. That is the accelerometer at work!
\end{experiment}

\sectionrule

% --- Gyroscope ---
\section*{Sensor 2: The Gyroscope}

The \term{gyroscope} measures how the phone is \emph{rotating} —
not just tilting, but spinning around any axis.
The accelerometer notices movement, while the gyroscope measures rotation in fine detail.

\begin{itemize}
  \item Games use the gyroscope so you can steer by tilting the phone.
  \item Panorama photo apps use it to stitch a wide photo accurately.
  \item Virtual reality apps use it to track which direction you are looking.
\end{itemize}

\begin{picturethis}
  Imagine spinning a toy top on a table.
  A gyroscope inside the phone works like a tiny, invisible top —
  it resists changes in direction, letting the phone measure exactly
  how it has turned.
\end{picturethis}

\sectionrule

% --- Light sensor ---
\section*{Sensor 3: The Light Sensor}

The \term{light sensor} (also called the ambient light sensor) measures
how bright the surroundings are.

\begin{itemize}
  \item On a sunny day, it tells the screen to get brighter so you can read it.
  \item In a dark room, it dims the screen to save your eyes and battery.
  \item It helps automatic brightness keep the screen comfortable to look at.
\end{itemize}

Phones also have a \term{proximity sensor}, which helps turn the screen off
when the phone is close to your face during a call.

\sectionrule

% --- Microphone ---
\section*{Sensor 4: The Microphone}

The \term{microphone} is a sensor that converts sound waves into electrical signals.

\begin{itemize}
  \item It lets you make calls and record voice notes.
  \item It powers voice assistants that listen for your commands.
  \item It records audio during video recording.
\end{itemize}

Most smartphones have two or three microphones to capture sound from
different directions and reduce background noise.

\sectionrule

% --- Compass ---
\section*{Sensor 5: The Compass}

The \term{compass} sensor (called a magnetometer) detects the direction
of Earth's magnetic field — just like a traditional compass needle.

\begin{itemize}
  \item It tells your maps app which direction you are facing.
  \item It helps GPS navigation show the correct arrow direction.
  \item Augmented reality apps use it to overlay digital objects on the real world.
\end{itemize}

\sectionrule

% --- Summary diagram ---
\begin{diagram}[The five hidden sensors and what each one detects.]
  \begin{tikzpicture}[>=Stealth]
    \node[draw=TechPurple,fill=TechPurple!15,rounded corners=8pt,
          minimum width=2cm,minimum height=3.2cm,line width=1.2pt,
          font=\bfseries\small,align=center] (phone) at (0,0) {Smart\\phone};
    % Sensors around
    \node[draw=LeafGreen!70,fill=LeafGreen!8,rounded corners,minimum width=3cm,minimum height=0.7cm,font=\small,right=2.2cm of phone,yshift=1.6cm]  (acc) {Accelerometer: movement};
    \node[draw=NeonGreen,fill=NeonGreen!8,rounded corners,minimum width=3cm,minimum height=0.7cm,font=\small,right=2.2cm of phone,yshift=0.5cm]  (gyr) {Gyroscope: rotation};
    \node[draw=SunYellow!80!black,fill=SunYellow!8,rounded corners,minimum width=3cm,minimum height=0.7cm,font=\small,right=2.2cm of phone,yshift=-0.5cm] (lgt) {Light sensor: brightness};
    \node[draw=WarmOrange!70,fill=WarmOrange!8,rounded corners,minimum width=3cm,minimum height=0.7cm,font=\small,right=2.2cm of phone,yshift=-1.6cm] (mic) {Microphone: sound};
    \node[draw=SmartBlue!70,fill=SmartBlue!8,rounded corners,minimum width=3cm,minimum height=0.7cm,font=\small,left=2.2cm of phone,yshift=0cm]    (cmp) {Compass: direction};
    % Arrows
    \draw[TechPurple!60,->] (phone.east) -- (acc.west);
    \draw[TechPurple!60,->] (phone.east) -- (gyr.west);
    \draw[TechPurple!60,->] (phone.east) -- (lgt.west);
    \draw[TechPurple!60,->] (phone.east) -- (mic.west);
    \draw[TechPurple!60,->] (phone.west) -- (cmp.east);
  \end{tikzpicture}
\end{diagram}

\sectionrule

\begin{newwords}
  \begin{description}[labelwidth=3.8cm, leftmargin=4.0cm]
    \item[\term{Accelerometer}]  Sensor that measures speed and direction of movement.
    \item[\term{Gyroscope}]      Sensor that measures rotation.
    \item[\term{Light sensor}]   Sensor that measures ambient brightness.
    \item[\term{Microphone}]     Sensor that converts sound waves into electrical signals.
    \item[\term{Compass}]        Sensor that detects Earth's magnetic field to find direction.
    \item[\term{Sensor}]         A device that detects changes in the physical world.
  \end{description}
\end{newwords}

\sectionrule

\begin{bigidea}
  A smartphone has many \textbf{hidden sensors} that detect the physical world:
  movement, rotation, light, sound, and magnetic direction.
  They work quietly in the background to make the phone smarter.
\end{bigidea}

\character{Sensa}{We sensors feed information to Chip so the phone
  can react to the world. Next, let’s meet the programmes that put
  all that power to work. Meet Appy in Chapter 8!}
